\section{Experimental Methodology}

\subsection{Benchmark Matrix}
Microbenchmarks sweep a full factorial over:
\begin{itemize}
    \item Particle counts $N_p \in \{10^4, 10^5, 10^6\}$.
    \item Grid sizes $64^2$, $128^2$, and $256^2$ on the unit square.
    \item Schemes: NGP, CIC, TSC, Esirkepov.
    \item Layouts: AoS and SoA.
    \item Sorting: enabled and disabled.
\end{itemize}
Each case uses 5 timed repeats after 3 warmup depositions.
Particles are initialized from a uniform Maxwellian with fixed seed 42.
Esirkepov microbenchmarks use $\Delta t = 0.01$ to form trajectory weights.
CPU scaling cases fix $N_p=10^5$, grid $128^2$, CIC, SoA, no sorting, and thread counts $T \in \{1,2,4,8\}$ with privatized OpenMP buffers.
Energy-drift simulations run 500 steps for all schemes at $N_p=10^5$ on a $128^2$ grid with sorting enabled.

\subsection{Timing Protocol}
Timed deposition separates sort cost from kernel cost.
When sorting is enabled, one sort is timed, then the deposit kernel is averaged over \texttt{repeats} iterations on the sorted particle order.
Throughput is reported as $N_p / t_{\mathrm{deposit}}$ and effective bandwidth as $(B_{\mathrm{read}}+B_{\mathrm{write}})/t_{\mathrm{deposit}}$, where $B_{\mathrm{read}}$ counts particle record bytes and $B_{\mathrm{write}}$ counts stencil updates plus one grid clear.
Sort time is excluded from throughput and bandwidth so that locality effects can be compared independently.

\subsection{Physics Validation}
Langmuir-wave validation uses $N_p=200$ particles on a $64^2$ grid with $\Delta t=0.002$ and 16384 timesteps.
A sinusoidal velocity perturbation of amplitude 0.05 is imposed at mode number $m=1$.
After each deposition, a neutralizing background $\rho_{\mathrm{bg}} = 1/(L_x L_y)$ is added.
Field energy is recorded over the final three-quarters of the run and analyzed by discrete Fourier transform.
Pass criterion: measured frequency spread across all four schemes $\leq 5\%$.
We also report $\omega_{\mathrm{meas}}/\omega_{p,\mathrm{macro}}$ to document the systematic offset from the 1D cold-plasma estimate in this 2D diagnostic setup.

Energy-drift validation uses separate 500-step full-PIC simulations at $N_p=10^5$ on a $128^2$ grid.

\subsection{Hardware and Scope}
All benchmarks were built with \texttt{cmake -DCMAKE\_BUILD\_TYPE=Release} and executed on the platform summarized in Table~\ref{tab:hardware}.
GPU evaluation is deferred to future work on CUDA hardware; CUDA kernels compile with \texttt{BUILD\_CUDA=ON} but are not timed in this study.
All reported figures use deposit-kernel time unless explicitly labeled otherwise.

\begin{table}[t]
\centering
\small
\caption{Benchmark platform configuration.}
\label{tab:hardware}
\begin{tabular}{ll}
\toprule
Component & Specification \\
\midrule
CPU & Apple M1 Pro \\
Cores & 8 (OpenMP threads: 1, 2, 4, 8) \\
OS & macOS 26.5.1 \\
Compiler & Apple clang 21.0.0, C++17, \texttt{-O3} \\
Libraries & FFTW 3.3.11, libomp 22.1.7 \\
\bottomrule
\end{tabular}
\end{table}
